\section{Classical foundations used}\label{sec:foundations}

We state the classical inputs to make the scope of the proof explicit. We also use elementary linear algebra, field and group theory, and the basic definitions and functorial formalism of schemes, sheaves, divisors, and K\"ahler differentials. The counting, index, inseparable-genus, valued-field, and specialization assertions particular to this proof are established in the later sections. In particular, no weight-counting theorem or conjectural reduction is an input.

\begin{foundation}[Finite groups]\label{foundation:groups}
For a finite group $X$, its $p$-subgroups lie in Sylow subgroups of order $|X|_p$. A nontrivial finite $p$-group has nontrivial center and an element of order $p$, and each of its maximal proper subgroups is normal. We use the orbit--stabilizer formula and the usual fixed-point and conjugation actions.
\end{foundation}

\begin{foundation}[Finite-dimensional algebras and characters]\label{foundation:algebra}
The Jacobson radical of a finite-dimensional algebra is nilpotent. Over an algebraically closed field, its semisimple quotient is a product of full matrix algebras, one for each simple module up to isomorphism. A finite-dimensional commutative algebra over an algebraically closed field is a product of local algebras with that field as residue, indexed by its primitive idempotents. Finite group algebras in characteristic zero are semisimple. Over a splitting field the simple endomorphism algebras are the scalar field, and an ordinary irreducible character satisfies
\[
 \sum_{x\in X}\chi(x)\chi(x^{-1})=|X|.
\]
Idempotents lift uniquely from the residue algebra of a finite commutative algebra over a complete discrete valuation ring. Nakayama's lemma applies to finite modules over any local ring.
\end{foundation}

\begin{foundation}[Commutative algebra]\label{foundation:commalg}
A finite-type algebra over a Noetherian ring is Noetherian. Integral extensions satisfy lying over. Closed points of a variety over a field are dense and have finite residue fields over the ground field. For an integral variety, the local dimension at a point plus the transcendence degree of its residue field is the transcendence degree of the function field. A regular local ring is a normal domain; a part of a minimal system of generators of its maximal ideal is a regular sequence and its quotient is regular of the correspondingly smaller dimension. A one-dimensional Noetherian normal local domain is a discrete valuation ring, as is a Noetherian local domain with nonzero principal maximal ideal. Finite torsion-free modules over discrete valuation rings are free. Completion of a discrete valuation ring retains its residue field and uniformizer.
\end{foundation}

\begin{foundation}[Fields and absolute field differentials]\label{foundation:fields}
If a finite group $D$ acts faithfully by automorphisms on a field $M$, then $M/M^D$ is Galois of degree $|D|$. Finite separable field extensions are simple. A finite field extension has a maximal separable subextension, and the remaining extension is purely inseparable; the number of embeddings in an algebraic closure is the separable degree. Multiplicative groups of finite fields are cyclic. Algebraically closed fields are classified by their characteristic and transcendence degree. If $S$ is a finitely generated field of transcendence degree $\rho$ over a perfect field $\kappa$, then
\[
 \dim_S\Omega_{S/\kappa}=\rho;
 \qquad [S:S^p]=p^\rho\quad\text{in characteristic }p.
\]
\end{foundation}

\begin{foundation}[Valuations over a complete discrete base]\label{foundation:valuation}
A valuation extends to an algebraic field extension. A complete discrete valuation on a field extends uniquely to any finite field extension; its valuation ring is the integral closure of the original ring, is finite over it, and is again a complete discrete valuation ring. The fields to which we apply the latter assertion initially have characteristic zero. No corresponding finiteness assertion for an arbitrary nondiscrete valuation ring is assumed.
See \cite[Tag 00IA]{Stacks} for prolongation and
\cite[Tags 032W, 09EM and 0325]{Stacks} for the complete
discrete-base finiteness assertions.
\end{foundation}

\begin{foundation}[Descent and finite models]\label{foundation:descent}
A projective scheme with a semilinear descent datum under a finite Galois field extension descends if equipped with a compatible linearized ample line bundle. The same statement holds for a finite \emph{\'etale} faithfully flat Galois extension of rings and a relatively ample line bundle. Equivariant morphisms and invariant closed subschemes descend along with the scheme. Flatness, smoothness, finiteness, \emph{\'etaleness}, and relative ampleness in these settings are checked after the faithfully flat extension. Geometric integrality of a finite-type scheme over a field can be checked after an extension of that field. A finite diagram of projective schemes and morphisms of finite presentation over a directed union of fields is defined over a finite stage, including its maps to a base scheme already defined there and the identities among those maps. One may realize projective descent by descent of semilinear modules and graded algebras followed by $\Proj$.
For the projective and module descent assertions, see \cite[Tags 0CCJ, 0CDR and 0D1V]{Stacks}.
The five locality assertions are respectively
\cite[Tags 02L2, 02VL, 02LA, 02VN and 0D3C]{Stacks}; geometric
integrality is covered by \cite[Tags 038K, 0384 and 054P]{Stacks}.
The finite-stage assertion follows from \cite[Tags 01ZM and 01ZP]{Stacks},
including chosen projective embeddings among the finite data.
\end{foundation}

\begin{foundation}[Normalization]\label{foundation:normalization}
The normalization of an integral variety over a field in a finite extension of its function field is finite. The normalization of a finite-type $\mathbb Z$-domain in its fraction field is finite. In the integral closure of a normal domain in a finite Galois extension of its fraction field, the Galois group acts transitively on the primes above any fixed prime of the base.
The normalization finiteness assertions follow from the Nagata property in \cite[Tag 0335]{Stacks}.
\end{foundation}

\begin{foundation}[Regularity, smoothness, and properness]\label{foundation:regularity}
Finite separable scalar extension preserves regularity and reducedness. Smooth varieties over a field are regular. A regular local ring essentially of finite type over a perfect field $\kappa$ is smooth at the corresponding point of a finite-type $\kappa$-model; its module of absolute differentials is free of rank the transcendence degree of the function field over $\kappa$. We use the valuative criterion for properness, including the unique extension of maps from fraction fields of valuation rings, and that proper maps are closed.
For regularity and smoothness over a perfect field, see \cite[Tags 00TV and 0B8X]{Stacks}.
\end{foundation}

\begin{foundation}[Proper cohomology and ample section rings]\label{foundation:cohomology}
Over a Noetherian base, the higher direct images of coherent sheaves under a proper morphism are coherent. Over an affine base their cohomology commutes with flat affine base change. For a projective scheme over an affine Noetherian base and a relatively ample line bundle, sufficiently high powers give projective embeddings, and sufficiently high twists annihilate higher coherent cohomology. The ample section ring is finitely generated and its $\Proj$ recovers the scheme. Finite pullback preserves ampleness; tensor products of relatively ample line bundles are relatively ample, and a line bundle is relatively ample if a positive power is. Formation of $\Proj$ commutes with base change. On a projective curve over a field, coherent cohomology vanishes in degrees greater than one, and for sheaves of finite support in degrees greater than zero. Global functions on a geometrically integral projective variety over a field equal the ground field.
For proper coherence, flat cohomology base change, and ample twists, respectively, see \cite[Tags 02O5, 02KH and 0B5T]{Stacks}.
For Proj recovery and the ampleness rules, see
\cite[Tags 0EKI, 0892, 0890 and 02NN]{Stacks}.
\end{foundation}

\begin{foundation}[Differentials, adjunction, and duality]\label{foundation:duality}
We use the transitivity sequence
\[
 \Omega_{V/U}\otimes_VW\longrightarrow\Omega_{W/U}
 \longrightarrow\Omega_{W/V}\longrightarrow0
\]
and the conormal sequence for a closed immersion. Let $S$ be a field. On $\mathbb P^b_S$,
$\det\Omega_{\mathbb P^b_S/S}=\mathcal O(-b-1)$. For a coherent sheaf $\mathcal G$ there, Serre duality identifies
\[
 H^i(\mathcal G)^\vee
 \simeq\operatorname{Ext}^{b-i}
       (\mathcal G,\mathcal O(-b-1)).
\]
We use the local-to-global Ext spectral sequence and the Koszul resolution of a quotient by a regular sequence. These inputs will yield the required absolute determinant-degree formula even when the curve is not smooth over $S$.
For the projective-space dualizing complex and the adjunction formalism, see \cite[Tag 0A9W and Section~48.15]{Stacks}.
The local-to-global Ext sequence is \cite[Tag 0BQP]{Stacks}.
\end{foundation}

\begin{foundation}[Elliptic curves]\label{foundation:elliptic}
A smooth Weierstrass cubic with its point at infinity is an elliptic curve. A generalized Weierstrass equation over a discrete valuation ring with unit discriminant gives a smooth elliptic scheme, polarized by three times its origin. Its relative cotangent line is trivial with invariant differential $\eta$, and multiplication by an integer $a$ pulls $\eta$ back to $a\eta$. For a positive integer $a$, multiplication by $a$ on an elliptic curve or elliptic scheme is finite locally free of degree $a^2$.
For the relative group law, see
\cite[II, Proposition~2.7]{DeligneRapoport1973}; the unit-discriminant
Weierstrass family and the invariant cotangent description are given in
\cite[Tags 072S and 047I]{Stacks}. See also
\cite[Chapter~III]{Silverman2009} for the Weierstrass formulation.
Multiplication is finite flat and surjective by
\cite[Proposition~20.7]{MilneAV2022}; over the Noetherian bases here it
is finite locally free by \cite[Tag 02KB]{Stacks}.
\end{foundation}

\begin{foundation}[Complex curves and covers]\label{foundation:complex}
Finite \'etale algebraic covers of a smooth complex algebraic curve
correspond to finite unramified topological covers of its analytification,
including morphisms. For a punctured projective curve, normalization of
the original projective curve in the resulting algebraic function field
gives the finite compact cover. Normal complex curves are smooth, and the
local punctured covering completes by a power map, so monodromy cycle
lengths give ramification indices. Topological covers are classified by
fundamental-group actions on finite fibers. The fundamental group of a
genus-one surface punctured at $n$ points has generators
$x,y,u_1,\ldots,u_n$ with the single relation
\[
 [x,y]u_1\cdots u_n=1.
\]
For a finite map of smooth proper curves over an algebraically closed field of characteristic zero, Riemann--Hurwitz states
\[
 2g(C)-2=\deg(f)\bigl(2g(C')-2\bigr)
              +\sum_{P\in C}(e_P-1).
\]
For the algebraic/topological cover correspondence, see \cite[Theorem~3.4]{MilneLEC2008}; for Riemann--Hurwitz, see \cite[Section~53.12, Tag 0C1B]{Stacks}.
\end{foundation}

\begin{foundation}[Noetherian connectedness]\label{foundation:connectedness}
A proper morphism over a locally Noetherian base whose structure-sheaf direct image is canonically the structure sheaf of the base has geometrically connected fibers. This is the connected-fiber assertion of Stein factorization. We apply it only after constructing a Noetherian normal model, not directly over the nondiscrete valuation ring.
See \cite[Tag 03H0]{Stacks}.
\end{foundation}

\subsection{Coefficient fields and scalar extension}\label{sec:coefficients}

We justify the coefficient choice in \Cref{sec:introduction} without
requiring a mixed-characteristic lift of every algebraically closed
field. Put $\kappa=\overline{\mathbb F}_p$ and choose an embedding
$\kappa\hookrightarrow k$. For any finite group $X$, the ideal
$\rad(\kappa X)\otimes_\kappa k$ is nilpotent, and the quotient of
$kX$ by this ideal is the scalar extension of the split matrix product
$\kappa X/\rad(\kappa X)$. That quotient is again a split semisimple
algebra. The ideal is therefore exactly $\rad(kX)$, and the simple
modules correspond under scalar extension.

The center also commutes with this extension: it is the kernel of the
finitely many linear maps $a\mapsto ax-xa$, $x\in X$, and field
extension preserves kernels. Each local factor of $Z(\kappa X)$ has
nilpotent radical and residue $\kappa$; after extension it still has
nilpotent radical, now with residue $k$, so remains local. Thus
primitive blocks correspond and their simple-module counts agree.
Coefficient restriction, restriction to centralizers in the Brauer
map, and the quotient maps for normal subgroups all commute with
scalar extension. The residue characters of the local central factors
extend as well. Consequently the central-character block assignment
in \eqref{eq:block-assignment} is preserved, for every subgroup and
subquotient involved in a weight. It suffices to prove the numerical
identity over $\kappa$.

To construct the coefficient ring, start with $\mathbb Z_p$, lift a
tower of finite separable extensions of its residue field $\mathbb F_p$
that exhausts $\kappa$, and complete the union. The construction at the
start of the proof of \Cref{prop:mixed-lift} shows that this gives a
complete discrete valuation ring of characteristic zero with residue
$\kappa$.
In an algebraic closure of its fraction field, choose matrix
realizations of every ordinary irreducible representation of every
subgroup of $G$. There are only finitely many such representations and
matrix entries. Adjoining their entries gives a finite field extension
that splits all those subgroups. Its valuation ring is again a complete
DVR by \Cref{foundation:valuation}; its finite residue extension is
still $\kappa$, which is algebraically closed. This gives the system
$(K,\mathcal O,\kappa)$ used in \Cref{sec:group}. Quotient
representations inflate, so the same coefficients split the required
subquotients. These coefficient fields are separate from the auxiliary
geometric fields constructed later.

The remaining sections prove the blockwise chain identity, the marked
high-index twist, the genus estimate with exact rounding, the required
valued-field assertions, and the specialization model. The foundations
above supply their classical inputs, not the desired weight equality.
